\documentclass[11pt]{article}
\usepackage[letterpaper,margin=0.78in]{geometry}
\usepackage{amsmath,amssymb,amsthm,mathtools}
\usepackage{graphicx,microtype,enumitem,xcolor,hyperref}
\hypersetup{colorlinks=true,hypertexnames=false,linkcolor=blue!45!black,urlcolor=blue!55!black,citecolor=blue!45!black}
\setlist{nosep,leftmargin=1.6em}
\newtheorem{theorem}{Theorem}
\newtheorem{lemma}{Lemma}
\newtheorem{remark}{Remark}
\newcommand{\R}{\mathbb R}
\newcommand{\cE}{\exp L^p}
\newcommand{\norm}[1]{\left\lVert #1\right\rVert}
\newcommand{\ip}[2]{\left\langle #1,#2\right\rangle}
\definecolor{answer}{RGB}{62,75,145}

\title{Spectral-gap decay and the ground-state asymptotic\\
\large A substantial answer to the decay-range question in arXiv:2306.02828}
\author{Automated mathematical discovery run \texttt{fa\_banach\_001}, lane 9}
\date{17 August 2026}

\begin{document}
\maketitle

\begin{center}
\fcolorbox{answer}{answer!6}{\parbox{0.91\linewidth}{
\textbf{Candidate substantial partial result.}
For the small-data global solution constructed in the source paper, the
reported polynomial decay is never the true large-time scale.  After a
possibly smaller smallness threshold, the solution decays in every
\(L^a(\R^d)\), \(p\leq a\leq\infty\), and in \(\exp L^p\), at every
exponential rate below \(d^\beta\).  More precisely, in \(L^2\) it equals a
multiple of the harmonic-oscillator ground state times
\(e^{-d^\beta t}\), up to a remainder with every rate below
\(\min\{(d+2)^\beta,m d^\beta\}\).
}}
\end{center}

\section*{Source question and transcription}
The source studies
\[
 \partial_tu+H^\beta u=f(u),\qquad
 H=-\Delta+|x|^2,\qquad 0<\beta\leq1,
\]
with small initial data in \(\exp L^p(\R^d)\) and a nonlinearity satisfying
\begin{equation}
 |f(s)-f(t)|\leq C|s-t|
 \bigl(|s|^{m-1}e^{\lambda|s|^p}
      +|t|^{m-1}e^{\lambda|t|^p}\bigr).                 \tag{S}
\end{equation}
Its Theorem 1.2 constructs a global weak-mild solution with uniform Orlicz
control and a polynomial \(L^a\)-estimate for a restricted range of
exponents.  Remark 1.2(iii) asks, verbatim: ``A natural question arises as
to whether the range of \(a\) required for (1.5) is optimal.''

This packet completely determines the large-time mechanism, but does not
settle the source's separate short-time endpoint restrictions.

\section*{Original source image}
\begin{center}
\includegraphics[width=0.97\textwidth]{assets/source_question_crop.png}
\end{center}
\noindent\footnotesize Source: Bhimani--Majdoub--Manna,
arXiv:2306.02828v3, PDF p.~4, Remark 1.2.  This is a direct crop from
\texttt{source\_paper.pdf}.
\normalsize

\section*{Definitions and source input}
The Luxemburg norm on \(\cE(\R^d)\) is
\[
 \norm{v}_{\cE}
 =\inf\left\{A>0:
 \int_{\R^d}\left(e^{|v(x)|^p/A^p}-1\right)dx\leq1\right\}.
\]
The source fixed-point argument supplies at least one admissible exponent
\(a_0>p\), numbers \(M,\sigma_0>0\), and a global solution satisfying
\begin{equation}
 \sup_{t>0}\norm{u(t)}_{\cE}\leq M,
 \qquad \norm{u(t)}_{a_0}\leq M t^{-\sigma_0}.          \tag{1}
\end{equation}
The small-data threshold can be decreased so that \(M\) is as small as
needed below.  We use two estimates already recorded in the source:
\begin{align}
 \norm{v}_q&\leq C_{p,q}\norm{v}_{\cE},
       &&p\leq q<\infty,                                \tag{2}\\
 \norm{e^{-tH^\beta}g}_q
 &\leq C e^{-d^\beta t}\norm{g}_r,
       &&t\geq1,\quad 1\leq r,q\leq\infty,             \tag{3}\\
 \norm{e^{-tH^\beta}g}_q
 &\leq C t^{-\frac d{2\beta}|1/r-1/q|}\norm{g}_r,
       &&0<t\leq1.                                      \tag{4}
\end{align}
Here and below, endpoint estimates are available because \(0<\beta\leq1\).

\section*{Idea / proof intuition}
The restricted exponent \(a_0\) is only a bootstrap seed.  Its polynomial
decay, together with uniform Orlicz moments, makes \(f(u(t))\) decay in a
sufficiently high \(L^r\).  Restarting the mild equation one time unit in
the past then forces \(u(t)\to0\) in \(L^\infty\).  Once this happens,
(S) becomes
\(|f(u)|\leq\kappa\delta^{m-1}|u|\), where \(\delta\) can be made
arbitrarily small by moving the restart time.  The spectral gap in (3)
therefore dominates the nonlinearity and yields exponential decay in
\(L^p\cap L^\infty\).

Finally, the discrete spectrum
\(d^\beta,(d+2)^\beta,\ldots\) separates the ground state from all higher
modes.  The nonlinear forcing is of order \(e^{-m\omega t}\), so a scalar
variation-of-constants formula gives the leading ground-state coefficient,
while the orthogonal complement decays at the minimum of the next spectral
rate and the nonlinear rate.

\section*{Main result}
Let \(\phi_0(x)=\pi^{-d/4}e^{-|x|^2/2}\), normalized in \(L^2\), and put
\[
 \lambda_0=d^\beta,\qquad \lambda_1=(d+2)^\beta.
\]

\begin{theorem}[Large-time range and asymptotic profile]
Assume the hypotheses of the source's global-existence theorem:
\(1<p\leq d(m-1)/(2\beta)\), \(m\geq p\), \(0<\beta\leq1\), and (S),
with \(f(0)=0\).  After possibly decreasing its small-data constant, the
global solution constructed there has the following properties.

\begin{enumerate}[label=(\roman*),topsep=3pt,itemsep=2pt]
\item For every \(a\in[p,\infty]\) and every \(0<\omega<\lambda_0\),
there are \(T_\omega,C_{a,\omega}>0\) such that
\[
 \norm{u(t)}_a+\norm{u(t)}_{\cE}
 \leq C_{a,\omega}e^{-\omega t},\qquad t\geq T_\omega. \tag{5}
\]
\item There is a scalar \(c\) such that, for every
\(0<\mu<\min\{\lambda_1,m\lambda_0\}\),
\[
 \norm{u(t)-c e^{-\lambda_0t}\phi_0}_2
 \leq C_\mu e^{-\mu t}                                \tag{6}
\]
for all sufficiently large \(t\).
\end{enumerate}
\end{theorem}

\begin{proof}
We divide the argument into four steps.

\paragraph{Step 1: a decaying forcing norm.}
Choose
\[
 r>\max\left\{\frac d{2\beta},\frac{a_0}{m}\right\}.
\]
By decreasing the source smallness threshold, assume that the uniform bound
in (1) is small enough that
\(2r\lambda(2M)^p\leq1\).  The Luxemburg definition and
\((e^x-1)^{2r}\leq e^{2rx}-1\) give
\begin{equation}
 \sup_{t>0}\norm{e^{\lambda|u(t)|^p}-1}_{2r}\leq C.    \tag{7}
\end{equation}
The Orlicz embedding (2) uniformly bounds every finite \(L^q\)-norm.
Since \(a_0<mr<2mr<4mr\), interpolation between the decaying
\(L^{a_0}\)-norm and the bounded \(L^{4mr}\)-norm gives positive numbers
\(\theta_1,\theta_2\) such that
\begin{align*}
 \norm{u(t)}_{mr}&\leq C(1+t)^{-\theta_1\sigma_0},\\
 \norm{u(t)}_{2mr}&\leq C(1+t)^{-\theta_2\sigma_0}.
\end{align*}
Putting the second scalar argument equal to zero in (S) yields
\(|f(v)|\leq C|v|^m e^{\lambda|v|^p}\).  Hence, by (7),
\begin{align}
 \norm{f(u(t))}_r
 &\leq C\norm{u(t)}_{mr}^m
 +C\norm{u(t)}_{2mr}^m
       \norm{e^{\lambda|u(t)|^p}-1}_{2r}
 \leq C(1+t)^{-\eta}                                  \tag{8}
\end{align}
for some \(\eta>0\).

\paragraph{Step 2: entry into a small \(L^\infty\)-ball.}
For \(t\geq2\), restart the mild equation at \(t-1\):
\[
 u(t)=e^{-H^\beta}u(t-1)
 +\int_{t-1}^t e^{-(t-s)H^\beta}f(u(s))\,ds.           \tag{9}
\]
The fixed-time \(L^{a_0}\)-to-\(L^\infty\) estimate and (1) make the
first term tend to zero.  For the integral, (4), (8), and
\(\gamma=d/(2\beta r)<1\) give
\[
 \left\lVert\int_{t-1}^t e^{-(t-s)H^\beta}f(u(s))\,ds
 \right\rVert_\infty
 \leq C(1+t)^{-\eta}\int_0^1\tau^{-\gamma}\,d\tau
 \longrightarrow0.                                    \tag{10}
\]
Thus \(\norm{u(t)}_\infty\to0\).

\paragraph{Step 3: spectral-gap bootstrap.}
Fix \(0<\omega<\lambda_0\) and choose
\(\omega<\nu<\lambda_0\).  Combining the short- and long-time parts of
(3)--(4), for \(q=p,\infty\) there is \(A\) such that
\begin{equation}
 \norm{e^{-tH^\beta}}_{L^q\to L^q}\leq A e^{-\nu t},
 \qquad t\geq0.                                        \tag{11}
\end{equation}
Choose \(\delta>0\) so small that
\(A C e^{\lambda\delta^p}\delta^{m-1}<\nu-\omega\), and then choose
\(T\) so that \(\norm{u(t)}_\infty\leq\delta\) for \(t\geq T\).
Condition (S) gives pointwise
\[
 |f(u(t,x))|\leq C e^{\lambda\delta^p}\delta^{m-1}|u(t,x)|.
\]
Taking the sum norm on \(B=L^p\cap L^\infty\), (9) restarted at \(T\)
and (11) imply
\[
 \norm{u(t)}_B\leq A e^{-\nu(t-T)}\norm{u(T)}_B
 +A\kappa_\delta\int_T^t e^{-\nu(t-s)}\norm{u(s)}_B\,ds,
\]
where \(\kappa_\delta=C e^{\lambda\delta^p}\delta^{m-1}\).
Weighted Gronwall therefore gives
\begin{equation}
 \norm{u(t)}_p+\norm{u(t)}_\infty\leq C_\omega e^{-\omega t}.
                                                               \tag{12}
\end{equation}
Interpolation proves (5) for every \(p\leq a\leq\infty\).  The source
inclusion
\(L^p\cap L^\infty\hookrightarrow\exp L^p\), with
\(\norm{v}_{\cE}\leq C(\norm{v}_p+\norm{v}_\infty)\), proves the Orlicz
part of (5).

\paragraph{Step 4: the ground-state coefficient.}
At a sufficiently large positive time \(T_2\), the solution belongs to
\(L^2\).  Indeed, choose \(T_2\) so that the already proved
\(L^\infty\)-smallness holds on \([T_2-1,T_2]\).  In the one-unit restart
formula the linear term is in \(L^2\) by fixed-time smoothing from \(L^p\).
On the last unit interval, (S) and boundedness give
\(|f(u)|\leq C|u|^m\), while interpolation between \(L^p\) and
\(L^\infty\) gives \(u\in L^{2m}\); hence the forcing is in \(L^2\).
Thus (9) is an \(L^2\) identity from \(T_2\) onward.

Let \(P_0v=\ip{v}{\phi_0}\phi_0\) and \(Q=I-P_0\).  Fix
\(\mu<\min\{\lambda_1,m\lambda_0\}\), and choose
\(\alpha<\lambda_0\) so close to \(\lambda_0\) that
\(m\alpha>\max\{\mu,\lambda_0\}\).
By (12), interpolation, and \(2m\geq p\),
\[
 \norm{f(u(t))}_2
 \leq C\norm{u(t)}_{2m}^m
 \leq C e^{-m\alpha t}.                                \tag{13}
\]
The scalar ground-state equation gives
\[
 e^{\lambda_0t}\ip{u(t)}{\phi_0}
 =e^{\lambda_0T_2}\ip{u(T_2)}{\phi_0}
 +\int_{T_2}^t e^{\lambda_0s}\ip{f(u(s))}{\phi_0}\,ds.
\]
The integral converges by (13).  Define
\begin{equation}
 c=e^{\lambda_0T_2}\ip{u(T_2)}{\phi_0}
 +\int_{T_2}^\infty e^{\lambda_0s}
     \ip{f(u(s))}{\phi_0}\,ds.                         \tag{14}
\end{equation}
The omitted tail is \(O(e^{-m\alpha t})\) after restoring the factor
\(e^{-\lambda_0t}\).  On the orthogonal complement, the spectral theorem
and (13) yield
\[
 \norm{Qu(t)}_2
 \leq e^{-\lambda_1(t-T_2)}\norm{Qu(T_2)}_2
 +C\int_{T_2}^t e^{-\lambda_1(t-s)}e^{-m\alpha s}\,ds
 \leq C_\mu e^{-\mu t}.
\]
Together with the ground-state tail, this is (6).
\end{proof}

\section*{Consequences and sharpness}
If the coefficient \(c\) in (14) is nonzero, (6) shows that the exact
large-time exponent is \(d^\beta\); no faster uniform estimate can hold for
that solution.  If \(c=0\), the solution decays faster, with the next
possible obstruction coming from either the first excited spectral level
or the order-\(m\) nonlinear forcing.  Thus the source's polynomial bound
is useful at short time and as a bootstrap seed, but is far from optimal at
large time.

For example, when \(f\equiv0\) and \(u_0=\varepsilon\phi_0\), one has
exactly \(u(t)=\varepsilon e^{-d^\beta t}\phi_0\), so the spectral rate is
attained within the source hypotheses.

\section*{Verification, novelty, and scope}
\begin{itemize}
\item Every integrability exponent used in Step 1 is finite and larger than
the decaying seed exponent; the uniform high moments come directly from the
Orlicz bound.
\item The near-time singularity in Step 2 is integrable precisely because
\(r>d/(2\beta)\).
\item The perturbative constant in Step 3 tends to zero with
\(\norm{u(t)}_\infty\); this is why every rate below, rather than merely one
unspecified rate below, \(d^\beta\) is obtained.
\item The \(L^2\) asymptotic does not assume \(u_0\in L^2\); fixed-time
harmonic-oscillator smoothing supplies \(L^2\) membership.
\item Formula (14) explicitly identifies the leading coefficient and the
remainder estimate separates the next spectral and nonlinear rates.
\end{itemize}

\paragraph{Novelty audit.}
Cheap run indexes were searched for arXiv:2306.02828 and the terms
``optimal range'', ``spectral gap'', ``ground-state asymptotic'', and
``Orlicz exponential nonlinearity''.  Bounded arXiv/primary-source searches
through 17 August 2026 found work on exponential nonlinearities and recent
sharp decay estimates for other oscillator equations, but no source giving
this all-\(a\) large-time theorem or the coefficient formula (14) in the
present Orlicz setting.  Novelty confidence is moderate, not definitive.

\paragraph{Limitations.}
The theorem answers the large-time part of the source question and gives
the precise asymptotic scale.  It does not enlarge the source's admissible
range in its exact all-time estimate (1.5), whose power of \(t\) also
encodes the singular behavior as \(t\downarrow0\).  It applies to the
small-data solution constructed by the source after a possibly smaller
threshold, not to arbitrary large-data solutions.  The coefficient \(c\)
may vanish for special data.

\paragraph{Human review.}
Pending.  Priority checks are the passage from the source weak-mild
representative to the one-unit restart identity in Lebesgue spaces, and the
uniform exponential-moment choice after shrinking the source smallness
constant.

\section*{References}
\begin{enumerate}[label={[\arabic*]}]
\item D. G. Bhimani, M. Majdoub, and R. Manna,
\emph{Heat equations associated to harmonic oscillator with exponential
nonlinearity}, arXiv:2306.02828v3, especially Theorem 1.2, Remark 1.2, and
Theorem 2.1.
\item D. G. Bhimani, R. Manna, F. Nicola, S. Thangavelu, and
S. I. Trapasso, \emph{On heat equations associated with fractional harmonic
oscillators}, Fract. Calc. Appl. Anal. 26 (2023), 2470--2492.
\item M. Majdoub and S. Tayachi, \emph{Global existence and decay estimates
for the heat equation with exponential nonlinearity}, Funkcial. Ekvac. 64
(2021), 237--259.
\end{enumerate}

\end{document}
